\heading{Part 2: audit an analysis you did not write}{Summary \hfill 12 October 2026}
\textbf{An AI reads faster than you do, which makes it a useful auditor and an unreliable one.} Your job is to direct the audit and then to check it. Repeat the pattern: \textbf{ask $\rightarrow$ inspect $\rightarrow$ verify $\rightarrow$ iterate}, not prompt $\rightarrow$ trust $\rightarrow$ done.

\topic{Give every claim a status}
\renewcommand{\arraystretch}{1.12}
\begin{tabularx}{\linewidth}{@{}p{1.55in}X@{}}
\toprule
\textbf{Status} & \textbf{When to use it}\\
\midrule
Supported & The evidence and a cited reference agree with the claim.\\
Questionable & The claim may hold, but the wording or method is doubtful.\\
Incorrect & The evidence or a cited reference contradicts the claim.\\
Insufficient evidence & These files cannot show it either way. Not the same as wrong.\\
\bottomrule
\end{tabularx}

\topic{How to judge the audit itself}
\begin{itemize}
  \item Which problems did it find? Which did it miss? Did it report anything that was not a problem?
  \item Are its citations accurate? Open the section it names and read it.
  \item Could you check each explanation yourself? Verify at least three findings.
  \item What change to the prompt or the context improved the result?
  \item Which kinds of problems were easy for it, and which were hard?
\end{itemize}

\topic{Git and GitHub, in one box}
\renewcommand{\arraystretch}{1.1}
\begin{tabularx}{\linewidth}{@{}p{1.35in}X@{}}
\textbf{Git, GitHub} & Git keeps a project's history on your computer. GitHub hosts a copy for backup and sharing. \textbf{Push} sends commits there; \textbf{pull} brings changes back.\\
\textbf{Commit} & A labeled checkpoint of the project, with a message. Not the same as saving a file.\\
\textbf{Root} & The folder that contains \code{.git}. Git tracks the project relative to it.\\
\textbf{Never nest} & Do not put one Git repository inside another. The outer one does not track the inner one's files, and tools can disagree about which history applies.\\
\textbf{.gitignore} & Lists files Git should skip: caches, environments, and large or regenerable output.\\
\textbf{What goes where} & Code, scripts, text, Markdown, configuration, documentation, and small data files belong in GitHub. Large raw data, PDFs, images, and binaries belong in cloud or institutional storage, such as Google Drive, OneDrive, shared lab storage, or a research data repository. A README can say where things live. My working limit is about 10 MB per Git file. Notebooks are the awkward middle case.\\
\end{tabularx}

\topic{Reference texts to give Claude}
\begin{itemize}
  \item \href{https://openintro-ims.netlify.app/}{\emph{Introduction to Modern Statistics}} (open textbook), free PDF from \href{https://www.openintro.org/book/ims/}{OpenIntro}.
  \item \href{https://www.itl.nist.gov/div898/handbook/}{NIST/SEMATECH \emph{e-Handbook of Statistical Methods}}, Chapters 1 and 7.
  \item \href{https://doi.org/10.1080/00031305.2016.1154108}{The ASA statement on p-values} (2016), two pages of principles.
\end{itemize}
Save copies in \code{references/pdfs/}. Check every citation Claude gives by opening the section it names.

\vfill
\small\repourl\par
\textbf{Turn over for the activity.}
